\documentclass[12pt,a4paper]{article}
\usepackage[T2A]{fontenc}
\usepackage[utf8]{inputenc}
\usepackage[russian]{babel}
\usepackage{amsmath,amssymb}
\usepackage{geometry}
\geometry{margin=2cm}

\title{Расчетное задание №1}
\author{Шубин Егор Вячеславович P3209}
\date{}

\begin{document}
\maketitle

\section*{1. Электростатика: метод суперпозиции}

\subsection*{Задание 1}
Пусть ось цилиндра совпадает с осью $z$, основание лежит в плоскости $z=0$, а поле ищется в центре основания.
Линейная плотность заряда вдоль оси = $\lambda$, тогда заряд тонкого кольца толщины $dz$:
\[
dq=\lambda\,dz.
\]
Вклад кольца в осевую компоненту поля:
\[
dE_z=\frac{1}{4\pi\varepsilon_0}\frac{dq\,z}{(R^2+z^2)^{3/2}}
=\frac{\lambda}{4\pi\varepsilon_0}\frac{z\,dz}{(R^2+z^2)^{3/2}}.
\]
\[
E=\frac{\lambda}{4\pi\varepsilon_0}\int_0^\infty \frac{z\,dz}{(R^2+z^2)^{3/2}}
=\frac{\lambda}{4\pi\varepsilon_0R}.
\]
\subsection*{Задание 2}
$r=R$: $\sigma=kR\cos\theta$.

\textbf{а) Потенциал и напряженность в центре.}
\[
\varphi(0)=\frac{1}{4\pi\varepsilon_0R}\oint \sigma\,dS.
\]
Так как $\sigma\propto \cos\theta$, интеграл по сфере равен нулю:
\[
\varphi(0)=0
\]
Поле в центре:
\[
\mathbf E(0)= -\frac{1}{4\pi\varepsilon_0R^3}\oint \sigma\,\mathbf r'\,dS
=-\frac{kR}{3\varepsilon_0}\,\mathbf e_z.
\]
Следовательно,
\[
\mathbf E(0)=-\frac{kR}{3\varepsilon_0}\mathbf e_z.
\]

\textbf{б) Производные.}
Т.к. $\nabla\varphi=-\mathbf E$:
\[
\frac{\partial\varphi}{\partial x}=0,\qquad
\frac{\partial\varphi}{\partial y}=0,\qquad
\frac{\partial\varphi}{\partial z}=\frac{kR}{3\varepsilon_0}.
\]

\section*{2. Электростатика: теорема Гаусса}

\subsection*{Задание 3}
По закону Гаусса:
\[
E(r)\,4\pi r^2=\frac{Q_{\text{внутр}}(r)}{\varepsilon_0}.
\]
Для $r\le R$:
\[
Q(r)=4\pi\int_0^r \rho_0\!\left(1-\frac{r'}{R}\right)r'^2\,dr'
=4\pi\rho_0\!\left(\frac{r^3}{3}-\frac{r^4}{4R}\right),
\]
\[
E_{\text{внутр}}(r)=\frac{\rho_0}{\varepsilon_0}\left(\frac r3-\frac{r^2}{4R}\right).
\]
Полный заряд шара:
\[
Q_\Sigma=Q(R)=\frac{\pi\rho_0R^3}{3},
\]
для $r\ge R$:
\[
E_{\text{внеш}}(r)=\frac{\rho_0R^3}{12\varepsilon_0r^2}.
\]

\[
\frac{dE_{\text{in}}}{dr}=\frac{\rho_0}{\varepsilon_0}\left(\frac13-\frac{r}{2R}\right)=0
\Rightarrow r_m=\frac{2R}{3},
\]
\[
E_{\max}=E(r_m)=\frac{\rho_0R}{9\varepsilon_0}.
\]

\subsection*{Задание 4}
Для $r>R$:
\[
Q_{\text{внутр}}(r)=Q_0+4\pi\int_R^r \frac{\alpha}{r'}r'^2dr'
=Q_0+2\pi\alpha(r^2-R^2).
\]
\[
E(r)=\frac{Q_{\text{внутр}}(r)}{4\pi\varepsilon_0r^2}
=\frac{\alpha}{2\varepsilon_0}
\frac{1}{\varepsilon_0r^2}\left(\frac{Q_0}{4\pi}-\frac{\alpha R^2}{2}\right).
\]
Чтобы $E$ не зависело от $r$, коэффициент при $1/r^2$ должен быть нулем:
\[
Q_0=2\pi\alpha R^2.
\]
\[
E=\frac{\alpha}{2\varepsilon_0}.
\]
\newpage
\section*{3. Энергия электрического поля}

\subsection*{Задание 5}
Пусть $\varepsilon(x)=\varepsilon_1+\dfrac{\varepsilon_2-\varepsilon_1}{d}x$, $x\in[0,d]$.
 $D=q/S$, а
\[
E(x)=\frac{D}{\varepsilon_0\varepsilon(x)}.
\]
\[
U=\int_0^d E(x)\,dx=\frac{q}{S\varepsilon_0}\int_0^d\frac{dx}{\varepsilon(x)}
=\frac{q}{S\varepsilon_0}\frac{d}{\varepsilon_2-\varepsilon_1}\ln\frac{\varepsilon_2}{\varepsilon_1}.
\]
Отсюда емкость:
\[
C=\frac{q}{U}=
\frac{S\varepsilon_0(\varepsilon_2-\varepsilon_1)}
{d\ln(\varepsilon_2/\varepsilon_1)}.
\]

Поляризация:
\[
P=D\!\left(1-\frac1\varepsilon\right),
\qquad
\rho_{\text{св}}=-\nabla\!\cdot\!\mathbf P=-\frac{dP}{dx}
=-D\frac{1}{\varepsilon^2}\frac{d\varepsilon}{dx}.
\]
Так как $D=q/S$ и $\dfrac{d\varepsilon}{dx}=\dfrac{\varepsilon_2-\varepsilon_1}{d}$:
\[
\rho_{\text{св}}(\varepsilon)=
-\frac{q}{S}\frac{\varepsilon_2-\varepsilon_1}{d\,\varepsilon^2}.
\]

\subsection*{Задание 6}
Для тонкой сферической оболочки радиуса $r$:
\[
W=\frac{q^2}{8\pi\varepsilon_0\varepsilon r}.
\]
При $\varepsilon=1$:
\[
W=\frac{q^2}{8\pi\varepsilon_0 r}.
\]
\[
W_{r\to R}=W\!\left(1-\frac{r}{R}\right),
\]
поэтому доля энергии внутри сферы радиуса $R$:
\[
\eta=1-\frac{r}{R}.
\]
Для $r=1\ \text{см}$, $R=1\ \text{м}$:
\[
\eta=1-0.01=0.99.
\]
При $q=1\cdot10^{-19}\ \text{Кл}$ и $r=1\ \text{см}$:
\[
W\approx 4.5\cdot10^{-27}\ \text{Дж}.
\]

\section*{4. Постоянный ток}

\subsection*{Задание 7}
\[
G=\frac{1}{\rho}\frac{S}{d}.
\]
Из формулы емкости $C=\varepsilon_0\varepsilon\,S/d$:
\[
\frac{S}{d}=\frac{C}{\varepsilon_0\varepsilon},
\qquad
I=UG=\frac{UC}{\rho\varepsilon_0\varepsilon}.
\]
\[
U=2000\ \text{В},\ 
C=3000\ \text{пФ}=3.0\cdot10^{-9}\ \text{Ф},\ 
\rho=100\ \text{ГОм}\cdot\text{м}=10^{11}\ \Omega\cdot\text{м},\ 
\varepsilon=7.
\]
\[
I\approx 9.7\cdot10^{-7}\ \text{А}\approx 0.97\ \mu\text{А}.
\]

\subsection*{Задание 8}
Дано: $l=1000\ \text{м}$, $S=1\ \text{мм}^2=10^{-6}\ \text{м}^2$, $I=4.5\ \text{А}$.
\[
n=\frac{\rho_mN_A}{M}\approx
\frac{8.9\cdot10^3\cdot 6.02\cdot10^{23}}{63.5\cdot10^{-3}}
\approx 8.4\cdot10^{28}\ \text{м}^{-3}.
\]
\[
v=\frac{I}{neS}\approx 3.3\cdot10^{-4}\ \text{м/с}.
\]
\[
t=\frac{l}{v}\approx 3\cdot10^6\ \text{с}
\]
\[
N=nSl,\qquad E=\rho_{\text{Cu}}\frac{I}{S},\qquad F_\Sigma=NeE.
\]
\[
F_\Sigma=n\,e\,\rho_{\text{Cu}}\,I\,l\approx 1.0\cdot10^6\ \text{Н}.
\]

\end{document}
